\begin{frame}[t]{Critérios de correção e síntese}
\begin{itemize}
\item A regra define o resultado esperado sem ambiguidade.
\item O pseudocódigo respeita a ordem e trata os dois ramos.
\item O fluxograma representa a mesma solução do texto.
\item Os testes de mesa cobrem o valor limite.
\item A escolha do paradigma orienta a organização do programa.
\end{itemize}
\begin{block}{Próxima aula: tipos de dados e variáveis}
Como representar temperatura, quantidade de peças e estado de um botão na memória do programa?
\end{block}
\end{frame}
